\chapter{The two ends of the gauge: what moves a leukocyte}\label{ch:leukocyte}

\section{Why blood first, and neutrophils first}

The instrument reads blood first because it can: leukocytes are the majority of the DNA in a blood draw, their purified methylomes are among the best
measured of any human cells, and DNA mixtures of known composition let the composition step be checked against truth. Within blood it reads neutrophils
first, because they are usually the largest share of the DNA, so their reading carries the most signal (Chapter~\ref{ch:separation}). Everything in this
chapter beyond the neutrophil readings already reported is a hypothesis with a named test. None is yet a finding.

\section{The two ends are different physics}

The distinction is sharp: immune cells working hard to make repairs are not cells whose informational fidelity is compromised,
and an elevated reading is not the same event as a severely suppressed one.

\begin{description}
\item[Below Normal ($A<0.95$).] The cell holds less entropy on its identity sites than the healthy cell of its kind: fewer errors than healthy. On this
      reading the writing process has slowed or stopped and identity maintenance is failing: senescence, exhaustion, a lineage running down. Senescent
      fibroblasts read below Normal on their unmethylated channel (Chapter~\ref{ch:gauge}).
\item[Above Normal ($A>1.05$).] More entropy than the healthy cell of its kind. In an image, the printer is in overdrive: more is being written than
      the surface can hold as a pattern, and the surface saturates. This is the direction of the horizon analogy of Chapter~\ref{ch:floorbreach}.
\item[Activity is neither.] A neutrophil clearing an infection does more work while keeping its identity. The physics does not predict that a busy immune cell
      moves on the gauge, and nothing on the report may say that it does.
\end{description}
All three are \conjecture.

% COMMISSIONING-RETURN fig:p4_fraction: figure taken out 2026-10-03 (plotted development values from the removed fitted tare). Put back, redrawn from the commissioned chain; see development/COMMISSIONING_RETURNS.md.


\section{Four mechanisms, four tests}

\begin{center}\small
\begin{longtable}{@{}p{0.195\textwidth}p{0.127\textwidth}p{0.264\textwidth}p{0.313\textwidth}@{}}
\toprule mechanism & end & which cells & test, and where it stands \\\midrule\endhead
the leukocyte is the tumor: leukemias, lymphoma, myeloma; clonal hematopoiesis as the pre-malignant form & above & the affected lineage; the others near 1.00 & the affected lineage read against its own healthy floor. Not yet possible: on the neutrophil sites leukemic blood stayed Normal in 6 of 10, and progenitor floors on EPIC are not built \\
aging & above & slowly, per person, over years & serial readings of one person; serial mode is not yet built (Chapter~\ref{ch:serial}) \\
T-cell exhaustion, immunosenescence & below & CD4 and CD8 specifically, with neutrophils near 1.00 in the same specimen & on single molecules, IAM-A now reads CD8 T cells against their own reference (commissioned on healthy CD8 T cells); on arrays, T-cell floors are not yet built \\
emergency hematopoiesis: immature cells pushed out of the marrow & composition first & immature fraction up; immature cells may read below Normal & needs a progenitor group in Stage~A \\\bottomrule
\end{longtable}
\end{center}

\section{Aging}

The direction of aging on the neutrophil gauge is not yet measured by the commissioned chain. \openprob{} One development reading on 450K arrays, outside the chain's platform, found
neutrophil Met-A rising with age among controls with $\rho=0.23$; the same arrays carried a slide effect large
enough to split the reading by sex, so the trend is not separated from the instrument. \observed{} One reading is that an old immune system has had more written
over a lifetime onto the same surface, with less headroom, which is the same direction as clonal hematopoiesis. \conjecture{} Whether an elevation in old blood is a slope
shared by every lineage or a clone riding on it is what per-cell trajectories in one person separate: a slope moves every lineage together; a clone moves one lineage away from
the others.

\section{What the leukocyte gauge does not do}

It does not detect a solid tumor through the blood cells' $A$. Immune cells are not expected to take on a tumor's state by fighting it. \conjecture{} A solid cancer, if it shows in whole blood at all, shows through composition or through shed cells. Chain v4 reads composition against IAM-Atlas v2.1, whose 79 cell types include the cells of solid organs (Chapter~\ref{ch:atlas}, Section~\ref{sec:atlas:v21story}); the Deconvolver is commissioned for the shares of four blood cell types, and its reading of shed organ cells in blood is a test still to run. Plasma cell-free DNA, where shed cells are a much larger share, and single molecules, where a tumor's held error is read directly, are the
routes for that question (Chapter~\ref{part4:ch:reach}).

\section{Why a reading of one cell does not cancel}

A score pooled over many sites and many cells, computed as a case group minus a control group, can move in opposite directions under different conditions without any cell going
both ways. Two things produce that. Pooling: some sites drift up and others down, and one lineage moves one way while another moves the other, netted by a composition that itself
moves with disease and age. And the comparison: a control group sitting above or below 1.00 flips the sign of a departure. Neither enters the chain. One cell is read on its own
identity sites, against its own healthy reference, and the methylated and unmethylated channels can be reported apart (Chapter~\ref{ch:gauge}).

\section{A count, not a score}

For conditions that involve the immune system as a whole, the intended report is how many immune cell types read outside Normal, and which: a count of cell readings, not one score
for the immune system as a whole. It needs a commissioned reading for each cell type and is not yet built. \openprob{} The test plan's immune track asks the same question organ by organ: whether the immune cells of whole blood carry a distinct reading for the organ in which a cancer is forming, read with Group-A on patients and healthy people run in the same series.
